\chapter{Introdução}
\label{1-introducao}

O presente capítulo visa apresentar o cenário no qual este trabalho está inserido, por meio da \hyperref[1-contextualizacao]{Contextualização}.
Bem como expor os \hyperref[1-questao]{Questionamentos} a serem respondidos com sua realização;
a \hyperref[1-justificativa]{Justificativa} para a realização deste trabalho;
o \hyperref[1-objetivos-gerais]{Objetivo Geral} e os \hyperref[1-objetivos-especificos]{Objetivos Específicos} a serem alcançados;
a \hyperref[1-metodologia]{Metodologia} que será aplicada.
Por fim, demonstrar os detalhes sobre a \hyperref[1-organizacao]{Organização} da monografia.

\section{Contextualização}
\label{1-contextualizacao}
% Apresente o contexto geral do problema
% Descreva a área de pesquisa e sua relevância atual

% 1. Apresentação da Doença e do Vetor
A dengue é uma infecção viral causada pelo vírus da dengue (DENV), transmitida aos humanos pela picada de mosquitos infectados. A doença é mais prevalente em regiões de clima tropical e subtropical, como no Brasil, onde as espécies do mosquito \textit{Aedes} (\textit{Aedes albopictus} e, especialmente, o \textit{Aedes aegypti}) atuam como os principais vetores \cite{ministerio2026dengue}. 

% Após se alimentar de um hospedeiro infectado, a fêmea do mosquito passa a ser vetora do vírus pelo restante de sua vida, uma vez que o agente patogênico replica-se em seu intestino e migra para as glândulas salivares. O intervalo entre a ingestão do vírus pelo mosquito e sua capacidade de transmiti-lo a um novo hospedeiro é denominado Período de Incubação Extrínseca (PIE). Esse período dura aproximadamente de 8 a 12 dias em temperaturas entre 25 °C e 28 °C, podendo variar conforme a amplitude das flutuações térmicas diárias, o genótipo viral e a concentração viral inicial \cite{who2024dengue}.

% 2. O Cenário Epidemiológico Crítico no Brasil
Historicamente, o Brasil apresenta uma situação epidemiológica da dengue considerada 
crítica. De acordo com o informe epidemiológico da Organização Pan-Americana da Saúde, 
\citeonline{opas2026dengue}, em 2025 o país registrou 1.655.644 casos suspeitos, com 
1.793 óbitos e uma taxa de incidência de 775,7 por 100.000 habitantes, configurando o 
maior volume absoluto de casos nas Américas.
Um fator agravante no país, dada a sua extensão territorial, é a circulação simultânea 
dos quatro sorotipos do vírus (DENV-1, DENV-2, DENV-3 e DENV-4), o que eleva 
substancialmente o risco de formas graves da doença, visto que infecções secundárias 
por sorotipos distintos estão intimamente associadas a casos graves.

Fatores sistêmicos como a urbanização acelerada e não planejada, a intensa mobilidade 
humana, as deficiências de saneamento básico e o armazenamento inadequado de água 
favorecem a proliferação do vetor. Somam-se a isso as mudanças climáticas globais, que 
tendem a expandir as condições favoráveis à reprodução e à atividade do 
\textit{Aedes aegypti}, potencialmente alcançando novas regiões geográficas 
\cite{who2024dengue}.

% 3. Vacinação vs. Controle Vetorial
%Em resposta a esse desafio, o Brasil tornou-se o primeiro país do mundo a disponibilizar vacinas contra a dengue no sistema público de saúde. A vacina Qdenga, indicada para indivíduos que ainda não contraíram a doença, demonstrou ser eficaz contra o DENV-1 em 69,8\% dos casos, contra o DENV-2 em 95,1\% e contra o DENV-3 em 48,9\%, com eficácia geral comprovada de 84,1\% na prevenção de hospitalizações \cite{agenciagov2024qdenga}. Embora a incorporação do imunizante ao Sistema Único de Saúde (SUS) em 2023 e ao Calendário Nacional de Vacinação em 2024 represente um avanço inquestionável, alcançando 521 dos 5.569 municípios brasileiros em sua primeira campanhas, o controle do vetor \textit{Aedes aegypti} permanece como principal estratégia de prevenção e controle da doença e de outras arboviroses urbanas como chikungunya e Zika \cite{ministerio2026dengue}.

% 4. A Limitação do Controle Tradicional e a Necessidade de Antecipação
% A complexidade do controle vetorial, contudo, é acentuada por um aspecto estrutural da dinâmica epidêmica: quando a transmissão já está estabelecida em uma região, com um ou mais sorotipos em circulação e elevada densidade de mosquitos, as ações reativas de controle mostram-se pouco efetivas. A velocidade de circulação viral supera amplamente o tempo necessário para a redução das populações do vetor. Nesses cenários, a queda no número de casos tende a ocorrer mais em função da imunidade de grupo adquirida pela população do que pela eficácia das intervenções aplicadas \cite{ministerio2026dengue}. Tal constatação reforça a importância crítica da atuação preventiva, tornando essencial o desenvolvimento de novas metodologias que apoiem a gestão pública antes da instauração de uma epidemia.

% 5. A Solução Tecnológica: A Essência e Importância dos Sistemas Multiagentes
Nesse contexto, a simulação computacional baseada em Sistemas Multiagentes (SMA), 
também conhecida como Modelagem Baseada em Agentes (MBA), desponta como uma abordagem 
tecnológica fundamental. Essa técnica permite representar computacionalmente entidades 
individuais do mundo real, tais como mosquitos, seres humanos, criadouros e fatores 
ambientais, dotando-os de comportamentos autônomos e regras de interação 
\cite{pascoe2023agent}.

A relevância dos SMA na epidemiologia reside na sua capacidade 
de mapear o fenômeno denominado ``micro-macro'': observar como interações locais e 
microscópicas geram padrões complexos de disseminação espacial em escala populacional. Para além de fornecer visualizações, os Sistemas Multiagentes operam como laboratórios 
virtuais seguros para análises do tipo \textit{what-if} (``e se?''). Isso possibilita 
a pesquisadores e gestores avaliar o comportamento do ecossistema e a eficácia de 
diferentes cenários de intervenção antes da aplicação de políticas no mundo físico 
\cite{pascoe2023agent}.

% 6. O Big Data e o Desafio da Integração de Dados Reais

No campo da modelagem computacional, a era do \textit{Big Data} transformou 
o potencial dessas simulações. Historicamente, modelos epidemiológicos baseavam-se em 
populações sintéticas e médias teóricas. Em um contexto de abundância informacional, 
porém, os dados passaram a assumir papel central nos sistemas computacionais e de apoio 
à decisão \cite{moreira2025data}.

% 7. A Subutilização e a Necessidade de Tratamento
Apesar desse imenso potencial, a simples existência dessa vasta quantidade de dados 
na saúde pública esbarra em um desafio prático recorrente: a subutilização para fins 
de gestão. Bases nacionais abertas, como o Sistema de Informação de Agravos de 
Notificação (SINAN), o Levantamento Rápido do Índice de Infestação por 
\textit{Aedes aegypti} (LIRAa) e o SISPNCD, agregam valor indispensável à modelagem. 
Não obstante, embora essenciais para o registro das notificações, esses sistemas 
historicamente operam de forma isolada e não integrada \cite{rizzi2016sigdengue}.

% 8. A Ancoragem da Simulação (Conclusão do raciocínio)
Diante dessas limitações, iniciativas que promovem o tratamento e a integração dessas 
bases de dados mostram-se um passo fundamental na aproximação entre o modelo 
computacional e a realidade epidemiológica. Nesse sentido, destaca-se o SIGDENGUE 
\cite{rizzi2016sigdengue}, desenvolvido pela UNIOESTE em parceria com a Prefeitura 
Municipal de Cascavel/PR, que integra as bases nacionais abertas e as complementa com 
registros locais do município, consolidando informações que historicamente eram 
acompanhadas de forma manual e isolada, e oferecendo suporte à tomada de decisão 
pelos gestores de saúde pública. O grupo responsável pelo SIGDENGUE também conduziu 
esforços paralelos de simulação com agentes computacionais \cite{rizzi2011computational}, 
evidenciando o potencial da integração entre dados reais e modelagem baseada em agentes 
no contexto de Cascavel.


% 9. O Escopo do Trabalho e Visão de Futuro (Fechamento)

Com base no exposto, o presente trabalho visa projetar e desenvolver um Sistema 
Multiagentes voltado à simulação da dinâmica de focos endêmicos da dengue. O escopo 
primário da pesquisa concentra-se na modelagem de agentes comportamentais cujas ações 
e interações no ambiente virtual sejam orientadas por dados reais, englobando variáveis 
georreferenciadas, climáticas e registros oficiais de saúde pública. Tal abordagem 
permite calibrar os agentes virtuais com comportamentos autênticos e condições 
geográficas precisas, conferindo ao sistema o papel de ferramenta de apoio estratégico 
aos gestores de saúde pública.

A arquitetura do simulador será concebida de forma modular, decisão que, embora não 
implique entregas adicionais no escopo deste trabalho, visa facilitar evoluções 
incrementais em pesquisas futuras, com vistas à disponibilização do sistema como 
software livre e de código aberto, fomentando a pesquisa na área de sistemas 
multiagentes e oferecendo à sociedade uma ferramenta acessível de apoio ao 
enfrentamento da dengue.

% Exemplo de como usar uma citação:
%\begin{citacao}
%A engenharia de software tem evoluído significativamente nos últimos anos \cite{inmetro2003}.
%\end{citacao}

\section{Questionamentos}
\label{1-questao}
% Liste aqui as questões que seu trabalho pretende responder
% Formule perguntas claras e objetivas
% As questões devem estar alinhadas com seus objetivos

Tendo em mente que a vigilância epidemiológica lida com um contexto de alta quantidade de dados governamentais — oriundos de várias fontes como SINAN e plataformas climáticas, muitas vezes subutilizados e não integrados —, a identificação ágil de surtos de dengue torna-se um desafio. Esse desafio é evidente no que tange à necessidade de cruzar variáveis ambientais complexas e à implementação de soluções computacionais capazes de modelar de forma realista o comportamento e a proliferação do vetor (\textit{Aedes aegypti}).

Assim sendo, este trabalho busca responder aos seguintes questionamentos:

\begin{enumerate}
    \item \textbf{Questão 1:} Como viabilizar a estruturação e a integração de dados abertos governamentais (epidemiológicos, climáticos e geográficos) para que atuem como parâmetros de entrada em um ambiente de simulação computacional?
    
    \item \textbf{Questão 2:} Como traduzir as regras biológicas de proliferação do \textit{Aedes aegypti} e suas interações com as flutuações climáticas em especificações comportamentais compatíveis com a arquitetura de um Sistema Multiagentes?
    
    \item \textbf{Questão 3:} Como a implementação desse Sistema Multiagentes, quando alimentado por dados empíricos, pode apoiar a identificação e a simulação de cenários de surtos de dengue em apoio à gestão pública?
\end{enumerate}

Com base no andamento do estudo exploratório realizado até o momento, pretende-se utilizar a modelagem de Sistemas Multiagentes (SMA) como ferramenta para representar o ecossistema estudado. Nessa modelagem, serão especificadas as relações de interdependência entre os agentes autônomos (como os mosquitos e os focos de água) e o ambiente, balizadas pelas variáveis meteorológicas e espaciais.

Para testar o modelo e a integração das bases de dados, será idealizado um cenário de simulação prático no qual a arquitetura se aplica, bem como implementada a execução do simulador utilizando recortes temporais e geográficos específicos. Entende-se por simulação, neste escopo, um recurso útil para demonstrar a dinâmica de um problema de saúde pública orientando-se pelas necessidades de análise espacial e identificação de padrões por parte dos gestores. Para facilitar a documentação e o entendimento dos resultados providos ao longo deste trabalho, serão realizadas execuções de validação comparando o comportamento do sistema com os dados históricos reais, conforme apresentado no Capítulo 4 - Metodologia.
\section{Justificativa}
\label{1-justificativa}
% Explique por que seu trabalho é importante
% Apresente argumentos que justifiquem a pesquisa

Computacionalmente, os Sistemas Multiagentes (SMA) constituem um paradigma baseado em entidades intrinsecamente distribuídas e dotadas de expressiva capacidade dedutiva. Diferente das arquiteturas monolíticas tradicionais, os agentes em um SMA caracterizam-se pela sua autonomia e proatividade, sendo capazes de interpretar percepções do ambiente e executar ações fundamentadas em regras de decisão \cite{wooldridge2009introduction}. Essa abordagem descentralizada permite a modelagem de indivíduos que colaboram ou competem dentro de um ecossistema virtual, tornando essa tecnologia ideal para representar domínios onde o conhecimento, as variáveis e o controle estão espacial e logicamente dispersos.

No contexto da epidemiologia, a identificação e o monitoramento de focos de dengue configuram um desafio analítico de alta complexidade que se alinha perfeitamente a essa arquitetura. Tais tarefas demandam a assimilação de insumos heterogêneos, que podem basear-se em densas séries históricas governamentais, em aquisições em tempo real de plataformas climáticas ou em levantamentos amostrais de campo. Tem-se, portanto, um processo operado por entradas massivas e constantes de dados. Para que essas informações sejam processadas de forma eficiente, é imperativa a distribuição de responsabilidades sistêmicas. Adicionalmente, há a necessidade premente de interpretação dessas entradas, acompanhada de deduções lógicas pontuais, tomadas de decisão estruturadas e ações de mitigação rápida \cite{moreira2025data}.

Para apoiar os gestores públicos nessas ações, as simulações despontam como recursos inestimáveis de planejamento estratégico. Elas atuam como laboratórios virtuais seguros, evitando ou mitigando problemas que podem ser identificados e testados previamente em ambiente de software. Surtos de arboviroses, como ocorre com a dengue, exemplificam problemas de saúde altamente dinâmicos e difíceis de prever por métodos lineares. Nesse aspecto, as simulações possuem um apelo visual intrínseco que facilita a percepção, a avaliação e o esclarecimento de situações reais complexas. A literatura corrobora o poder dessa abordagem metodológica em diversos domínios não determinísticos: em \citeonline{ref_exemplo_transito}, os autores utilizam simulações para criar percepções preditivas sobre a formação de gargalos no trânsito urbano; de forma análoga, em \citeonline{ref_exemplo_epidemia}, simula-se a dinâmica de aglomerações humanas e contágio, o que ajuda sobremaneira as autoridades a compreender o impacto espacial da dispersão de agentes patogênicos.

Diante do exposto, a aplicação de Sistemas Multiagentes no contexto do monitoramento do \textit{Aedes aegypti} torna-se altamente pertinente. O uso de SMA fornece o arcabouço computacional robusto para gerenciar a distribuição e a complexidade das variáveis epidemiológicas. Agregar a simulação a essa solução proposta vai além da simples análise de dados estáticos; permite criar um ambiente de experimentação contínua. Esse "casamento" tecnológico transforma grandes volumes de dados governamentais em inteligência acionável, oferecendo aos tomadores de decisão uma ferramenta preditiva que não apenas mapeia o passado, mas projeta cenários fundamentais para a prevenção e contenção efetiva da doença.


% Exemplo de citação que reforce seus argumentos:
%\begin{citacao}
%O desenvolvimento contínuo de novas tecnologias demanda constante atualização das práticas de engenharia \cite{Sommerville2007}.
%\end{citacao}

\section{Objetivos}
\label{1-objetivos}
% Apresente uma breve introdução aos objetivos do trabalho
Para responder às questões de pesquisa levantadas e solucionar a problemática da integração de dados reais na previsão epidemiológica, este trabalho define metas claras de execução. A seguir, apresentam-se o objetivo central da pesquisa e os objetivos específicos metodológicos necessários para alcançá-lo.

\subsection{Objetivo Geral}
\label{1-objetivos-gerais}

O objetivo geral deste trabalho é projetar e desenvolver um Sistema Multiagentes parametrizável focado na simulação espacial e comportamental da proliferação do \textit{Aedes aegypti}, construindo uma arquitetura de software modular capaz de integrar dados abertos georreferenciados, climáticos e de saúde pública para apoiar a identificação preditiva de surtos de dengue em diferentes regiões.

\subsection{Objetivos Específicos}
\label{1-objetivos-especificos}
Para viabilizar o alcance do objetivo geral, o escopo do projeto foi estruturado em objetivos específicos, categorizados de acordo com suas naturezas teórica, prática e documental:

\begin{itemize}
    \item \textbf{Objetivos Específicos de viés Teórico:}
    \begin{enumerate}
        \item \textbf{Investigar} os conceitos fundamentais de interesse deste trabalho, com destaque para: Saúde Pública e Dados Abertos, Arboviroses, Sistemas Multiagentes, Simulação, Usabilidade, além do levantamento de dados climáticos e geográficos (utilizando o Distrito Federal como cenário de referência inicial).
        \item \textbf{Identificar} e mapear as bases de dados abertos de saúde pública, climáticos e geográficos adequadas para atuarem como insumos e parâmetros na calibração da simulação almejada.
        \item \textbf{Compreender}, a partir da revisão da literatura especializada, a biologia e a dinâmica que regem o ciclo de proliferação do \textit{Aedes aegypti}, mapeando as variáveis que o influenciam.
    \end{enumerate}
    
    \item \textbf{Objetivos Específicos de viés Prático:}
    \begin{itemize}
        \item \textbf{Especificar} o processo de proliferação do vetor biológico utilizando a abordagem arquitetural de Sistemas Multiagentes Comportamentais, definindo formalmente as regras do ambiente e dos agentes.
        \item \textbf{Implementar} o Sistema Multiagentes Comportamental, traduzindo o modelo especificado em um simulador executável focado no processo de proliferação do \textit{Aedes aegypti} para a identificação preditiva de surtos de dengue.
        \item \textbf{Analisar} a solução desenvolvida de forma quantitativa e qualitativa, aplicando matrizes de avaliação para mapear suas Forças, Fraquezas, Oportunidades e Ameaças (análise SWOT).
    \end{itemize}
    
    \item \textbf{Objetivos Específicos de viés Documental:}
    \begin{itemize}
        \item \textbf{Documentar} integralmente a arquitetura estruturada, as especificações e os resultados obtidos, garantindo que a solução seja reprodutível, transparente e acessível à comunidade acadêmica e aos interessados em saúde pública.
    \end{itemize}
\end{itemize}

\section{Metodologia}
\label{1-metodologia}
% Apresente um resumo da metodologia utilizada
% Indique o tipo de pesquisa e principais métodos
% Detalhe completo será feito no capítulo de metodologia

Este trabalho enquadra-se como uma pesquisa de natureza aplicada e com objetivo exploratório. Quanto à abordagem do problema, adota-se um modelo de pesquisa híbrido (métodos mistos), contemplando análises quantitativas e qualitativas. Em relação aos procedimentos técnicos, a condução divide-se em bibliográfica, exploratória e experimental. Detalhes mais aprofundados sobre esses aspectos metodológicos serão expostos no Capítulo 4 - Metodologia.

Conferindo apenas um breve resumo sobre os métodos utilizados para o alcance das metas, tem-se: pesquisa bibliográfica para fundamentar conceitos teóricos e revisar modelos existentes, visando o cumprimento dos objetivos de viés teórico; e pesquisa experimental para o desenvolvimento, validação e teste do simulador proposto, cobrindo os objetivos de viés prático. Para o controle e a execução dessas atividades de construção de software, far-se-á o uso de um \textit{Backlog} estruturado acompanhado por um quadro Kanban \cite{chave_referencia_agil}, organizando as demandas de forma iterativa. Por fim, emprega-se a pesquisa documental para o cumprimento do objetivo de viés documental.

Cabe mencionar ainda que o modelo construído passará por um rigoroso processo de validação em duas frentes, refletindo a abordagem híbrida adotada. Na dimensão quantitativa, serão mensuradas as métricas de desempenho do sistema e a aderência da simulação às séries históricas reais, comprovando sua eficácia na identificação de surtos para subsidiar ações governamentais mais rápidas e direcionadas. Na dimensão qualitativa, será conduzida uma avaliação formal de usabilidade da ferramenta, verificando se o usuário final considera o sistema intuitivo no que tange à entrada de insumos, navegação e visualização de dados. Os resultados dessas avaliações serão devidamente analisados e reportados na documentação final.


\section{Organização da Monografia}
\label{1-organizacao}

Para uma melhor compreensão de como este documento está estruturado, a seguir é apresentada a organização da monografia, dividida em sete capítulos essenciais:

\begin{itemize}
    \item \textbf{Capítulo \ref{cap-referencial-teorico} -- \nameref{cap-referencial-teorico}:} Detalhamento sobre as principais referências teóricas utilizadas nesta pesquisa, abrangendo conceitos sobre Saúde Pública e Dados Abertos, Arboviroses e a biologia do \textit{Aedes aegypti}, além dos fundamentos de Simulação Computacional e Sistemas Multiagentes.

    \item \textbf{Capítulo \ref{cap-suporte-tecnologico} -- \nameref{cap-suporte-tecnologico}:} Exposição das ferramentas tecnológicas, linguagens de programação, \textit{frameworks} de simulação e bibliotecas de processamento de dados empregados para viabilizar a arquitetura e a construção da solução como um todo.

    \item \textbf{Capítulo \ref{cap-metodologia} -- \nameref{cap-metodologia}:} Classificação metodológica desta pesquisa, bem como a descrição detalhada dos métodos, procedimentos e ferramentas de gestão ágil (como o acompanhamento de \textit{Backlog} via Kanban) utilizados para guiar o cumprimento dos objetivos práticos e teóricos.

    \item \textbf{Capítulo \ref{cap-proposta} -- \nameref{cap-proposta}:} Apresentação aprofundada da proposta deste trabalho. Neste capítulo, detalha-se o processo de modelagem dos agentes comportamentais, a estruturação do \textit{pipeline} de integração de dados empíricos e a arquitetura do simulador de proliferação do vetor.

    \item \textbf{Capítulo \ref{cap-resultados} -- \nameref{cap-resultados}:} Apresentação e discussão dos resultados obtidos. Engloba a avaliação quantitativa da simulação (aderência aos dados históricos da dengue) e qualitativa (avaliação de usabilidade da ferramenta e levantamento SWOT da solução proposta).

    \item \textbf{Capítulo \ref{cap-status-atual} -- \nameref{cap-status-atual}:} Considerações finais sobre o \textit{status} atual do trabalho, discutindo o panorama dos resultados da primeira etapa de desenvolvimento (TCC 1), retomando o contexto, os objetivos alcançados e delineando os passos e desafios metodológicos projetados para a continuidade da pesquisa.
\end{itemize}